% 第五章 重点板块深度分析（二）：中周期拐点型

\chapter{重点板块深度分析（二）：中周期拐点型}

中周期拐点型板块的核心特征是：行业经历了1--2年的调整或蓄势，当前正处于景气度向上拐点的左侧或右侧初期，市场对其持续性仍存分歧，但产业层面的验证信号已逐步明确。这类板块的预期差来源于\textbf{市场对周期拐点的时间判断滞后于产业实际变化}，一旦拐点确认，估值与盈利有望形成戴维斯双击。

本章重点分析两个中周期拐点型板块：智驾链零部件（L3/L4渗透率拐点）和半导体设备材料（国产替代深化拐点）。两者均处于产业趋势明确但市场认知不充分的阶段，具备较高的赔率与胜率平衡。

\begin{keyinsight}[中周期拐点型板块核心特征]
\textbf{拐点左侧布局，戴维斯双击可期。}中周期拐点型板块的投资价值在于：市场认知滞后于产业实际变化，当拐点确认时，盈利预期上修与估值中枢提升形成共振。此类板块通常具备“业绩验证信号已现、市场共识尚未形成”的特征，是超额收益的重要来源。
\end{keyinsight}

\section{智驾链零部件：L3/L4渗透率拐点在即，产业链价值重构}

智驾链零部件是汽车电动化之后下一个十年维度的产业大趋势。随着高阶智驾从“尝鲜配置”向“大众标配”渗透，产业链价值正在从传统机械件向电子件、从硬件向软件、从分散部件向域集中架构快速迁移。我们判断，2026--2027年是L3级自动驾驶从“政策准入”走向“规模普及”的关键窗口期，产业链相关公司将迎来业绩与估值的双重重估。

\subsection{景气度三维验证}

\textbf{业绩端验证：营收加速增长，盈利能力持续改善}

2026年Q1，A股智驾链核心公司（域控制器、线控底盘、传感器、HUD等）合计营收同比增长27.3\%，较2025年全年的18.6\%显著加速；归母净利润同比增长34.8\%，利润增速高于收入增速，反映规模效应逐步显现。分环节看，域控制器环节增速最高（同比+42\%），线控底盘次之（+31\%），智能座舱（+24\%）与传感器（+22\%）稳健增长。

值得注意的是，\textbf{头部公司的增长速度显著快于行业平均}，呈现出明显的集中化趋势。德赛西威2026Q1智能驾驶业务收入同比增长58\%，新签订单同比增长65\%；伯特利线控制动产品营收同比增长76\%，渗透率快速提升。这表明智驾产业链正从“主题投资”阶段进入“业绩兑现”阶段。

\textbf{产业端验证：高阶智驾车型密集上市，渗透率快速跃升}

2026年以来，搭载L2+及以上智驾功能的新车型上市节奏明显加快。据乘联会统计，2026年1--5月，国内新上市乘用车中，搭载L2级及以上智驾功能的车型占比已达58\%，较2025年全年的42\%提升16个百分点；其中\textbf{搭载城市NOA功能的L2++车型占比从2025年的12\%跃升至28\%}，渗透率提升速度超市场预期。

\begin{exhibitbox}[图表5-1：中国乘用车智驾渗透率趋势图（2022--2027E）]
\centering
\vspace{0.3cm}
\begin{tikzpicture}
\begin{axis}[
    width=0.85\textwidth, height=6cm,
    ylabel={渗透率（\%）},
    xlabel={年份},
    ymin=0, ymax=80,
    xmin=2022, xmax=2027,
    xtick={2022,2023,2024,2025,2026,2027},
    xticklabels={2022,2023,2024,2025,2026E,2027E},
    legend pos=north west,
    legend style={font=\scriptsize},
    grid=major,
    grid style={dashed, rulegrey!50},
    tick label style={font=\scriptsize},
    label style={font=\small},
]
% L2及以上
\addplot[thick, navy, mark=circle, mark size=2.5pt] coordinates {
    (2022, 28) (2023, 35) (2024, 42) (2025, 48) (2026, 58) (2027, 70)
};
\addlegendentry{L2及以上}

% L2+/NOA
\addplot[thick, accentblue, mark=square, mark size=2.5pt] coordinates {
    (2022, 5) (2023, 9) (2024, 15) (2025, 22) (2026, 38) (2027, 55)
};
\addlegendentry{L2+/城市NOA}

% L3及以上
\addplot[thick, nvgreen, mark=triangle, mark size=3pt] coordinates {
    (2022, 0.2) (2023, 0.5) (2024, 1.2) (2025, 2.5) (2026, 8) (2027, 20)
};
\addlegendentry{L3及以上}

\end{axis}
\end{tikzpicture}
\vspace{0.2cm}
\sourcenote{乘联会、工信部、AStock研究团队测算。2026--2027年为预测值。}
\end{exhibitbox}

核心芯片出货量数据进一步验证了这一趋势。英伟达2026财年Q1（对应自然年Q1）汽车业务营收同比增长124\%，其中Orin/Xavier系列智驾芯片出货量同比增长132\%。国内方面，华为MDC系列智驾计算平台2026年Q1出货量同比增长超过200\%，地平线征程系列芯片出货量同比增长87\%。\textbf{芯片端的高增长是智驾渗透率提升最直接、最领先的指标}，其数据强度远超市场一致预期。

\textbf{政策端验证：L3准入政策落地，法规体系逐步完善}

2025年11月，工信部等四部委联合发布《关于开展智能网联汽车准入和上路通行试点的通知》，正式启动L3级智能网联汽车准入试点。2026年3月，首批L3准入车型名单公布，包括宝马i7、奔驰E级、智己L7等7款车型获得准入许可，标志着\textbf{L3级自动驾驶在中国正式进入合法上路阶段}。

地方层面，深圳、上海、北京、广州等城市已先后开放L3级自动驾驶公开道路测试，部分城市开始探索L4级无人出租、无人配送的商业化运营。政策面的持续加码为智驾产业发展提供了清晰的制度框架，也为消费者接受度提升奠定了基础。

\subsection{预期差来源分析}

我们认为，当前市场对智驾链零部件板块的认知存在三大预期差：

\textbf{预期差一：低估智驾Tier1国产替代的速度与深度}

市场普遍认为，智能驾驶Tier1市场长期被博世、大陆、采埃孚等海外巨头垄断，国产替代需要漫长过程。但实际情况是，\textbf{以华为、小鹏、理想为代表的中国造车新势力正在通过全栈自研+本土供应链的模式，快速打破海外Tier1的垄断格局}。

以域控制器为例，2023年海外厂商市占率仍超过80\%，到2025年已降至约55\%，预计2026年将进一步降至45\%以下。德赛西威、华为、地平线、知行科技等本土厂商的市占率快速提升。国产替代的核心驱动因素不仅是成本优势（本土方案价格通常比海外低20--30\%），更重要的是\textbf{响应速度和迭代能力}——本土Tier1能够更好地配合中国车企快速的车型迭代节奏，这是海外巨头难以匹敌的。

\textbf{预期差二：对单车价值量提升的幅度预期不足}

市场倾向于将智驾零部件公司按照传统汽车零部件的估值框架定价（15--20x PE），但忽视了智能驾驶带来的单车价值量呈指数级提升。传统燃油车中，汽车电子单车价值量约2,000--3,000元；到L2级时代提升至8,000--12,000元；\textbf{L3级时代单车智驾相关硬件价值量可达2.5--4万元，L4级更将突破5万元}，增长空间在10倍以上。

\begin{exhibitbox}[图表5-2：不同级别自动驾驶对应单车智驾价值量对比]
\centering
\vspace{0.3cm}
\begin{tikzpicture}
\begin{axis}[
    width=0.85\textwidth, height=6cm,
    ylabel={单车价值量（万元）},
    xlabel={自动驾驶级别},
    ymin=0, ymax=6,
    symbolic x coords={L0/L1, L2, L2+, L3, L4},
    xtick=data,
    tick label style={font=\scriptsize},
    label style={font=\small},
    bar width=0.6cm,
    nodes near coords,
    nodes near coords style={font=\scriptsize, /pgf/number format/.cd,fixed,precision=1},
]
\addplot[ybar, fill=navy!60, draw=navy] coordinates {
    (L0/L1, 0.3)
    (L2, 1.0)
    (L2+, 1.8)
    (L3, 3.2)
    (L4, 5.2)
};
\end{axis}
\end{tikzpicture}
\vspace{0.2cm}
\sourcenote{Yole、高工智能汽车、AStock研究团队整理。含感知、决策、执行全环节硬件价值。}
\end{exhibitbox}

更重要的是，智驾零部件的\textbf{软件价值占比持续提升}，域控制器、操作系统、算法等软件环节的利润率显著高于传统硬件。一旦软件收费模式跑通（如订阅制），相关公司的盈利模式将从“一次性硬件销售”转向“持续性服务收入”，估值中枢有望系统性上移。

\textbf{预期差三：忽视“智驾平权”带来的下沉市场增量}

市场普遍将高阶智驾定位为高端车型的配置，认为渗透率提升速度会受限于价格。但实际情况是，\textbf{智驾技术的下沉速度远超预期}——通过算法优化和硬件方案简化，L2++级智驾正在向15--20万元价位的车型渗透，甚至部分10--15万元车型也开始标配基础智驾功能。

华为的“乾崑”智驾方案、比亚迪的“天神之眼”、小鹏的“XNGP”等都在推出面向中端市场的减配版本，推动智驾从“豪华车专属”向“国民级配置”转变。这意味着智驾渗透率的天花板和提升速度都可能显著高于当前市场预期。

\subsection{核心催化与时间窗口}

我们认为，未来6--12个月智驾链有三大核心催化，有望推动板块估值系统性提升：

\textbf{催化一：L3级自动驾驶全国性推广（时间窗口：2026Q3--Q4）}

首批L3准入试点车型已在部分城市上路，预计2026年下半年将扩大试点范围至更多城市和更多车型。如果L3准入从“试点”转向“常态化”，将是行业的标志性事件，意味着智驾产业从“技术验证期”进入“商业化普及期”。

\textbf{催化二：城市NOA大规模推送（时间窗口：2026年全年持续）}

2026年被称为“城市NOA元年”，华为、小鹏、理想、特斯拉等头部企业都在加速推送城市NOA功能。预计到2026年底，\textbf{开通城市NOA功能的车辆保有量将从2025年底的约150万辆增长至500万辆以上}，覆盖城市数量从100+扩展到全国主要城市。城市NOA的大规模落地将直接验证高阶智驾的用户价值，推动消费者购车决策向智驾配置倾斜。

\textbf{催化三：华为/特斯拉智驾方案外溢（时间窗口：2026Q2--Q3）}

华为智驾方案正在从“问界独享”向“全行业开放”转型，奇瑞、江淮、北汽、长安等多个品牌的华为智驾车型将在2026年密集上市。特斯拉FSD入华进程也在加速推进，若FSD中国版正式落地，将对整个智驾产业生态产生深远影响。\textbf{头部企业的技术外溢将带动整个产业链的成熟度提升和成本下降}，惠及全产业链优质供应商。

\subsection{产业链图谱与受益环节}

智驾产业链可分为感知层、决策层、执行层三大环节，以及配套的软件与服务。

\begin{exhibitbox}[图表5-3：智驾产业链图谱与价值分布]
\centering
\vspace{0.3cm}
\resizebox{0.85\textwidth}{!}{\begin{tikzpicture}[
  layer/.style={rectangle, draw=navy, thick, fill=navy!15, text width=4.5cm, align=center, minimum height=1.2cm, font=\small\bfseries, rounded corners=3pt},
  sub/.style={rectangle, draw=steelgrey, fill=white, text width=2.3cm, align=center, minimum height=1.0cm, font=\scriptsize, rounded corners=2pt},
  subw/.style={rectangle, draw=steelgrey, fill=white, text width=2.8cm, align=center, minimum height=1.0cm, font=\scriptsize, rounded corners=2pt}
]

% --- 三层标题节点（先绘制，箭头在子节点下方穿过） ---
\node[layer] (perception) at (0, 4.5) {感知层（价值占比 25--30\%）};
\node[layer, fill=navy!25] (decision) at (0, 0.5) {决策层（价值占比 40--45\%）};
\node[layer] (execution) at (0, -3.5) {执行层（价值占比 25--30\%）};

% --- 层间箭头（先于子节点绘制，使子节点覆盖于箭头之上） ---
\draw[->, thick, navy] (perception.south) -- (decision.north);
\draw[->, thick, navy] (decision.south) -- (execution.north);

% --- 感知层子模块（4个，y=3.0 水平均匀分布，总宽约11.3cm） ---
\node[sub] (cam) at (-4.5, 3.0) {摄像头\\ \tiny 舜宇光学等};
\node[sub] (radar) at (-1.5, 3.0) {毫米波雷达\\ \tiny 保隆科技等};
\node[sub] (lidar) at (1.5, 3.0) {激光雷达\\ \tiny 禾赛/速腾等};
\node[sub] (ultra) at (4.5, 3.0) {超声波雷达\\ \tiny 奥迪威等};

% --- 决策层子模块（3个，y=-1.0 水平均匀分布，总宽约11.3cm） ---
\node[subw] (chip) at (-4.25, -1.0) {智驾芯片\\ \tiny 英伟达/华为/地平线};
\node[subw] (domain) at (0, -1.0) {域控制器\\ \tiny 德赛西威/华阳集团};
\node[subw] (sw) at (4.25, -1.0) {算法/软件\\ \tiny 中科创达等};

% --- 执行层子模块（3个，y=-5.0 水平均匀分布，总宽约11.3cm） ---
\node[subw] (brake) at (-4.25, -5.0) {线控制动\\ \tiny 伯特利/亚太股份};
\node[subw] (steer) at (0, -5.0) {线控转向\\ \tiny 耐世特/浙江世宝};
\node[subw] (sus) at (4.25, -5.0) {空气悬架\\ \tiny 保隆科技/中鼎股份};

\end{tikzpicture}}
\vspace{0.2cm}
\sourcenote{AStock研究团队整理。价值占比为L3级智驾单车硬件价值估算。}
\end{exhibitbox}

\begin{itemize}[leftmargin=2em]
\item \textbf{感知层}：摄像头、毫米波雷达、激光雷达、超声波雷达等。受益标的包括保隆科技（摄像头+毫米波雷达）、联合光电、舜宇光学等。该环节竞争相对充分，价值占比约25--30\%。
\item \textbf{决策层}：智驾域控制器、芯片、操作系统、算法等。\textbf{这是产业链价值量最高、技术壁垒最强的环节}，价值占比约40--45\%。受益标的包括德赛西威（域控制器龙头）、华阳集团（座舱域控+智能驾驶）、均胜电子（智驾域控+汽车电子整合）等。
\item \textbf{执行层}：线控制动、线控转向、空气悬架等。受益标的包括伯特利（线控制动龙头）、中鼎股份、保隆科技（空气悬架）等。该环节国产替代空间大，价值占比约25--30\%。
\end{itemize}

从投资优先级看，我们推荐\textbf{决策层优于执行层优于感知层}，主要原因：一是决策层技术壁垒最高，受益于智驾价值量提升的弹性最大；二是执行层受益于线控化趋势，国产替代确定性强；三是感知层竞争相对激烈，价格压力较大。

\subsection{风险因素与下行空间}

\textbf{主要风险因素：}

\begin{enumerate}[leftmargin=2em]
\item \textbf{智驾渗透率不及预期风险}：如果消费者对智驾功能的付费意愿不足，或者技术成熟度未达预期，可能导致渗透率提升慢于预期。下行空间：若渗透率增速放缓30\%，相关公司业绩增速可能下调20--25\%，对应估值承压。

\item \textbf{价格战侵蚀利润风险}：智驾行业竞争加剧，各车企可能通过降价来争夺市场份额，进而向供应链传导成本压力。尤其是域控制器、传感器等环节，价格下降速度可能快于预期。

\item \textbf{技术路线变革风险}：端到端大模型方案可能颠覆传统的“感知---决策---控制”分层架构，对现有Tier1厂商的技术路径形成冲击。如果特斯拉“无雷达纯视觉”方案或端到端大模型成为主流，部分硬件供应商可能面临价值重估。
\end{enumerate}

\textbf{下行空间评估}：当前智驾链核心标的2026E PE中位数约30倍，处于近三年50--60\%分位，估值已包含一定的成长预期。若行业景气度不及预期，估值可能回落至20--25倍区间，对应回调空间约20--30\%。但考虑到产业长期趋势明确，回调后的布局价值将更加显著。

\subsection{关键标的点评}

\textbf{德赛西威（002920）：智驾域控制器龙头，英伟达深度合作伙伴}

公司是国内智驾域控制器的绝对龙头，与英伟达、高通等芯片巨头深度合作，产品覆盖L2到L4各级别。2026Q1智能驾驶业务收入同比增长58\%，在手订单饱满。公司的域控制器产品已进入比亚迪、理想、小鹏、广汽等多家主流车企供应链，受益于高阶智驾渗透率提升的弹性最大。预计2026--2028年净利润复合增速约35\%，当前对应2026E PE约28倍。

\textbf{伯特利（603596）：线控制动龙头，底盘智化核心受益标的}

公司是国内线控制动（WCBS）的先行者和龙头，打破了博世等海外厂商的长期垄断。线控制动是L3级及以上自动驾驶的必备执行部件，渗透率将随着智驾等级提升快速提升。公司2026Q1线控制动业务收入同比增长76\%，新获定点项目持续增加。预计2026--2028年净利润复合增速约30\%，当前对应2026E PE约25倍。

\textbf{华阳集团（002906）：智能座舱+智驾域控双轮驱动}

公司在HUD（抬头显示）领域处于国内领先地位，同时布局座舱域控制器和智能驾驶域控制器，受益于座舱智能化与驾驶智能化双重趋势。HUD业务随车型升级持续放量，智驾域控业务与华为、地平线深度合作，增长潜力大。预计2026--2028年净利润复合增速约28\%，当前对应2026E PE约22倍。

\textbf{均胜电子（600699）：汽车电子平台型公司，智驾整合能力突出}

公司通过外延并购内生发展，已形成汽车安全、汽车电子两大业务板块，在智驾域控制器、智能座舱、车联网等领域均有布局。公司客户覆盖全球主流车企，受益于全球汽车智能化趋势。预计2026--2028年净利润复合增速约25\%，当前对应2026E PE约20倍。

\textbf{保隆科技（603197）：传感器+空气悬架双赛道龙头}

公司在汽车传感器领域品类齐全，摄像头、毫米波雷达、胎压监测等产品全面布局，同时空气悬架业务快速放量，受益于汽车智能化与消费升级双重逻辑。预计2026--2028年净利润复合增速约30\%，当前对应2026E PE约24倍。

\section{半导体设备材料：国产替代进入深水区，订单周期拐点临近}

半导体设备和材料是集成电路产业的基石，也是国产替代空间最大、战略意义最强的环节。在全球半导体周期触底回升、国内晶圆厂扩产持续、海外出口管制升级的多重背景下，半导体设备材料板块正处于\textbf{国产替代深化与周期复苏共振}的关键节点。我们判断，2026年下半年设备订单周期有望迎来明确拐点，材料验证加速也将带来业绩弹性。

\subsection{景气度三维验证}

\textbf{业绩端验证：头部公司维持高增长，盈利能力稳步提升}

2026年Q1，A股半导体设备板块（按申万分类）合计营收同比增长32.4\%，归母净利润同比增长41.8\%，连续第六个季度营收增速超过30\%。其中刻蚀设备、薄膜沉积设备、清洗设备等环节增速领先，分别达到38\%、42\%和35\%。半导体材料板块营收同比增长28.7\%，归母净利润同比增长36.2\%，增速较2025年下半年有所加快。

\textbf{公司层面的分化较为明显}：头部公司凭借技术壁垒和客户优势，业绩增长确定性强。北方华创2026Q1营收同比增长45\%，净利润同比增长53\%；中微公司营收同比增长38\%，净利润同比增长47\%；拓荆科技营收同比增长52\%，净利润同比增长68\%。而二三线设备公司增速相对平缓，反映行业正进入“强者恒强”的阶段。

\textbf{产业端验证：晶圆厂扩产持续，设备国产化率稳步提升}

国内晶圆厂扩产是半导体设备需求的核心支撑。2026年以来，中芯国际、华虹半导体、长江存储、长鑫存储等头部晶圆厂均公布了新的扩产计划。据SEMI统计，2026年中国大陆晶圆厂设备支出预计达到约230亿美元，同比增长18\%，占全球晶圆厂设备支出的比重约26\%，继续位居全球第一。

\begin{exhibitbox}[图表5-4：中国大陆晶圆厂设备支出及全球占比（2020--2027E）]
\centering
\vspace{0.3cm}
\begin{tikzpicture}
% 左轴：设备支出（柱状图）
\begin{axis}[
    width=0.85\textwidth, height=6cm,
    ylabel={设备支出（亿美元）},
    ymin=0, ymax=320,
    xmin=2019.5, xmax=2027.5,
    xtick={2020,2021,2022,2023,2024,2025,2026,2027},
    xticklabels={2020,2021,2022,2023,2024,2025,2026E,2027E},
    tick label style={font=\scriptsize},
    label style={font=\small},
    grid=major,
    grid style={dashed, rulegrey!50},
    legend pos=north west,
    legend style={font=\scriptsize},
    bar width=0.5cm,
]
% 中国大陆设备支出（柱状）
\addplot[ybar, fill=navy!60, draw=navy] coordinates {
    (2020, 125)
    (2021, 158)
    (2022, 185)
    (2023, 170)
    (2024, 195)
    (2025, 195)
    (2026, 230)
    (2027, 270)
};
\addlegendentry{中国大陆设备支出}

% 占比标注（柱顶文字）
\node[font=\tiny, color=nvgreen, above] at (axis cs:2020, 125) {占比：18\%};
\node[font=\tiny, color=nvgreen, above] at (axis cs:2021, 158) {占比：22\%};
\node[font=\tiny, color=nvgreen, above] at (axis cs:2022, 185) {占比：26\%};
\node[font=\tiny, color=nvgreen, above] at (axis cs:2023, 170) {占比：25\%};
\node[font=\tiny, color=nvgreen, above] at (axis cs:2024, 195) {占比：24\%};
\node[font=\tiny, color=nvgreen, above] at (axis cs:2025, 195) {占比：24\%};
\node[font=\tiny, color=nvgreen, above] at (axis cs:2026, 230) {占比：26\%};
\node[font=\tiny, color=nvgreen, above] at (axis cs:2027, 270) {占比：28\%};

\end{axis}

% 右轴：全球占比（折线图）
\begin{axis}[
    width=0.85\textwidth, height=6cm,
    ylabel={全球占比（\%）},
    ylabel style={color=nvgreen},
    yticklabel style={color=nvgreen},
    ymin=10, ymax=35,
    axis y line*=right,
    axis x line=none,
    xmin=2019.5, xmax=2027.5,
    xticklabels={,,},
    tick label style={font=\scriptsize},
    label style={font=\small},
    legend pos=north east,
    legend style={font=\scriptsize},
]
% 全球占比趋势线（实际百分比数值）
\addplot[thick, nvgreen, mark=o, mark size=2pt, mark options={fill=nvgreen}] coordinates {
    (2020, 18)
    (2021, 22)
    (2022, 26)
    (2023, 25)
    (2024, 24)
    (2025, 24)
    (2026, 26)
    (2027, 28)
};
\addlegendentry{全球占比趋势}

\end{axis}
\end{tikzpicture}
\vspace{0.2cm}
\sourcenote{SEMI、AStock研究团队测算。柱顶数字为全球占比（\%），2026--2027年为预测值。}
\end{exhibitbox}

更重要的是\textbf{设备国产化率的系统性提升}。2023年国内晶圆厂设备国产化率仅约15--20\%，到2025年已提升至约28--32\%，部分成熟制程产线国产化率已超过40\%。分设备类型看，刻蚀设备国产化率最高（约40\%+），薄膜沉积设备次之（约25--30\%），检测设备和光刻机国产化率仍然较低（不足10\%），但也在快速突破中。

材料端验证节奏稍慢但方向明确。硅片、特种气体、抛光液、靶材等关键材料的国产验证加速推进，部分材料已通过头部晶圆厂认证并实现批量供货。沪硅产业12英寸硅片2026Q1出货量同比增长62\%，安集科技抛光液在先进制程的验证取得积极进展。

\textbf{政策端验证：大基金三期落地在即，产业政策持续加码}

国家大基金三期设立的消息已在市场流传多时，据公开信息，三期基金规模有望达到3,000亿元以上，重点投向半导体设备、材料、先进制程等卡脖子领域。如果三期基金在2026年下半年正式落地，将显著提升市场对半导体设备材料板块的风险偏好，并为行业发展提供实质性资金支持。

此外，《集成电路产业高质量发展行动方案》、《“十四五”数字经济发展规划》等政策文件持续落地，从税收优惠、人才培养、应用牵引等多个维度支持半导体产业发展。\textbf{政策面的持续加持为国产替代提供了稳定的制度环境和资源保障}。

\subsection{预期差来源分析}

市场对半导体设备材料板块的分歧较大，主要存在三大预期差：

\textbf{预期差一：过度担忧周期下行，忽视国产替代的结构性机会}

市场普遍将半导体设备视为强周期行业，担心全球半导体周期下行会导致设备需求下滑。但我们认为，\textbf{中国大陆半导体设备市场具有显著的“逆周期”特征}——即使全球半导体行业处于下行周期，国内晶圆厂的扩产和国产替代进程也不会停止，反而可能因为海外设备交付受阻而加速国产替代。

从历史数据看，2022--2023年全球半导体行业经历了深度调整，全球半导体设备销售额同比下降约15\%，但中国大陆设备进口额和国产设备出货量仍保持正增长。这说明\textbf{国产替代逻辑在很大程度上对冲了周期下行的风险}。当前市场对周期下行的担忧可能过度定价了周期风险，而低估了国产替代的持续力度。

\textbf{预期差二：对先进制程突破的速度和意义认识不足}

市场普遍认为，国内半导体产业在先进制程（7nm及以下）上与海外差距巨大，短期难以突破。但实际上，国内在先进制程方面的进展快于市场预期。中芯国际的7nm工艺已实现小规模量产，长江存储的3D NAND技术已追近全球领先水平，刻蚀、薄膜沉积等关键设备在先进制程的验证也在稳步推进。

更重要的是，\textbf{先进制程突破带来的不仅是市场空间的打开，更是整个产业链信心的提升}。一旦某一款关键设备通过先进制程验证，其示范效应将带动其他设备和材料的验证提速，形成“以点带面”的突破格局。当前市场对先进制程突破的预期偏保守，一旦取得实质性进展，有望成为板块估值提升的重要催化剂。

\textbf{预期差三：低估半导体材料国产替代的加速度}

市场对半导体设备的国产替代关注度较高，但对半导体材料的国产替代进度关注不足。实际上，半导体材料的国产替代正在加速推进，部分领域的进度甚至快于设备。

与设备相比，半导体材料的验证周期更长、客户粘性更强，但一旦通过认证，供应商的市场地位非常稳固，盈利能力也更强。\textbf{当前国内半导体材料正处于“验证收获期”的前夜}——过去几年投入的大量研发和验证资源，将在未来2--3年逐步转化为订单和收入。市场尚未充分认识到这一转换期的业绩弹性。

\subsection{核心催化与时间窗口}

未来6--12个月，半导体设备材料板块有三大核心催化值得关注：

\textbf{催化一：大基金三期正式落地（时间窗口：2026Q3）}

大基金三期是市场高度关注的标志性事件。如果三期基金规模超预期（如达到3,500亿元以上），或投向比例向设备材料倾斜力度加大，将直接提升板块估值。更重要的是，大基金三期落地往往伴随着产业政策的系统性加码，可能配套税收优惠、政府采购等支持政策，形成“资金+政策”的双重催化。

\textbf{催化二：先进制程设备验证通过（时间窗口：2026Q3--Q4）}

国内多家设备厂商的刻蚀机、薄膜沉积设备、清洗设备等正在进行14nm甚至7nm制程的客户验证。如果下半年有设备正式通过先进制程验证并获得批量订单，将是国产替代进程中的标志性事件，有望显著提振市场信心。尤其是刻蚀和薄膜沉积两大核心设备品类的突破，将极大地打开国产设备的市场空间。

\textbf{催化三：出口管制倒逼替代（时间窗口：2026年全年持续）}

海外对中国半导体产业的出口管制呈现持续升级趋势。如果管制范围进一步扩大（如向更先进制程延伸、覆盖更多设备品类、限制人员交流等），短期内可能对市场情绪造成扰动，但中长期来看将\textbf{倒逼国内晶圆厂加速导入国产设备材料，缩短验证周期，提升采购比例}。

回顾历史，每一次出口管制升级之后，国产设备材料的导入速度都会明显加快。2022年10月美国对华先进计算芯片和半导体设备出口管制升级后，国内晶圆厂的设备国产化率提升速度从年均2--3个百分点加快到5--7个百分点。因此，出口管制升级短期是利空，中长期反而是国产替代的催化剂。

\subsection{产业链图谱与受益环节}

半导体设备产业链可分为前道设备、后道设备和核心零部件三大环节；半导体材料可分为制造材料和封装材料两大类。

\begin{exhibitbox}[图表5-5：半导体设备材料产业链图谱与国产化率]
\centering
\vspace{0.3cm}
\resizebox{0.9\textwidth}{!}{\begin{tikzpicture}[
  cat/.style={rectangle, draw=navy, thick, fill=navy!15, text width=3.2cm, align=center, minimum height=0.9cm, font=\small\bfseries, rounded corners=3pt},
  item/.style={rectangle, draw=steelgrey, fill=white, text width=3.2cm, align=center, minimum height=0.8cm, font=\scriptsize, rounded corners=2pt},
  sub/.style={rectangle, draw=steelgrey, text width=2.5cm, align=center, minimum height=1.1cm, font=\scriptsize, rounded corners=2pt}
]

% 图例（顶部居中）
\node[draw=steelgrey, fill=white, rounded corners=2pt, inner sep=5pt, font=\tiny] (legbox) at (0, 7.2) {
  \begin{tabular}{@{}c@{\hspace{3pt}}l@{\hspace{8pt}}c@{\hspace{3pt}}l@{\hspace{8pt}}c@{\hspace{3pt}}l@{}}
  \cellcolor{nvgreen!10}\rule[1.5pt]{0.35cm}{0.2cm} & 国产化率较高（$>$30\%） &
  \cellcolor{riskamber!15}\rule[1.5pt]{0.35cm}{0.2cm} & 国产化率中等（15--30\%） &
  \cellcolor{riskred!20}\rule[1.5pt]{0.35cm}{0.2cm} & 国产化率较低（$<$15\%）
  \end{tabular}
};

% 左列：设备
\node[cat] (equip) at (-6.5, 6.2) {半导体设备};
\node[item] (fequip) at (-6.5, 4.4) {前道设备（价值占比80--85\%）};
\node[item] (bequip) at (-6.5, 1.8) {后道设备（价值占比10--12\%）};
\node[item] (parts) at (-6.5, 0.3) {核心零部件（价值占比5--8\%）};

% 前道设备细分（右侧纵向排列）
\node[sub, fill=nvgreen!10] (etch) at (-2.8, 5.7) {\small 刻蚀\\\scriptsize 国产化率约40\%\\\scriptsize 中微/北方华创};
\node[sub, fill=riskamber!15] (depo) at (-2.8, 4.4) {\small 薄膜沉积\\\scriptsize 国产化率约28\%\\\scriptsize 拓荆/北方华创};
\node[sub, fill=riskred!20] (litho) at (-2.8, 3.1) {\small 光刻\\\scriptsize 国产化率约5\%\\\scriptsize 上海微电子};

% 后道设备细分
\node[sub, fill=nvgreen!10] (test) at (-2.8, 1.8) {\small 测试设备\\\scriptsize 国产化率约35\%\\\scriptsize 长川/华峰测控};

% 右列：材料
\node[cat] (material) at (6.5, 6.2) {半导体材料};
\node[item] (wmat) at (6.5, 3.5) {晶圆制造材料（价值占比70--75\%）};
\node[item] (pmat) at (6.5, 0.1) {封装材料（价值占比25--30\%）};

% 制造材料细分（左侧纵向排列）
\node[sub, fill=riskamber!15] (wafer) at (2.8, 5.4) {\small 硅片\\\scriptsize 国产化率约20\%\\\scriptsize 沪硅/立昂微};
\node[sub, fill=nvgreen!10] (gas) at (2.8, 4.1) {\small 特种气体\\\scriptsize 国产化率约30\%\\\scriptsize 华特/金宏};
\node[sub, fill=riskred!20] (photo) at (2.8, 2.8) {\small 光刻胶\\\scriptsize 国产化率约10\%\\\scriptsize 南大光电/晶瑞};
\node[sub, fill=riskamber!15] (cmp) at (2.8, 1.5) {\small CMP抛光材料\\\scriptsize 国产化率约15\%\\\scriptsize 安集/鼎龙};

% 底部标的栏
\node[draw=steelgrey, fill=panelgrey, rounded corners=2pt, text width=14cm, align=center, minimum height=0.9cm, font=\scriptsize] (targets) at (0, -1.2) {
  \textbf{核心受益标的}\quad 设备：中微公司、北方华创、拓荆科技、长川科技 \enskip\textperiodcentered\enskip 材料：沪硅产业、安集科技、华特气体、南大光电
};

% 连接线 - 设备层级（竖向相邻连接，避免箭头穿过中间盒子）
\draw[->, thick, navy] (equip.south) -- (fequip.north);
\draw[->, thick, navy] (fequip.south) -- (bequip.north);
\draw[->, thick, navy] (bequip.south) -- (parts.north);

% 连接线 - 前道设备 → 细分
\draw[->, thick, navy] (fequip.east) -- (etch.west);
\draw[->, thick, navy] (fequip.east) -- (depo.west);
\draw[->, thick, navy] (fequip.east) -- (litho.west);

% 连接线 - 后道设备 → 细分
\draw[->, thick, navy] (bequip.east) -- (test.west);

% 连接线 - 材料层级
\draw[->, thick, navy] (material.south) -- (wmat.north);
\draw[->, thick, navy] (wmat.south) -- (pmat.north);

% 连接线 - 制造材料 → 细分
\draw[->, thick, navy] (wmat.west) -- (wafer.east);
\draw[->, thick, navy] (wmat.west) -- (gas.east);
\draw[->, thick, navy] (wmat.west) -- (photo.east);
\draw[->, thick, navy] (wmat.west) -- (cmp.east);

\end{tikzpicture}}
\vspace{0.2cm}
\sourcenote{SEMI、AStock研究团队整理。国产化率为2025年底估算值。}
\end{exhibitbox}

\textbf{设备环节：}
\begin{itemize}[leftmargin=2em]
\item \textbf{前道设备}：刻蚀机、PVD/CVD/ALD等薄膜沉积设备、光刻机、涂胶显影设备、清洗设备、离子注入机、CMP设备等。这是价值量最高、国产替代空间最大的环节，价值占比约80--85\%。受益标的包括中微公司（刻蚀设备龙头）、北方华创（综合设备龙头）、拓荆科技（薄膜沉积设备）、芯源微（涂胶显影设备）、盛美上海（清洗设备）等。
\item \textbf{后道设备}：测试设备、分选机、探针台等。国产替代进度相对较快，部分领域已实现较高国产化率。受益标的包括长川科技、华峰测控等。
\item \textbf{核心零部件}：射频电源、真空泵、精密运动台、阀门等。这是设备产业链的“卡脖子”环节，国产化率最低，替代空间最大。
\end{itemize}

\textbf{材料环节：}
\begin{itemize}[leftmargin=2em]
\item \textbf{晶圆制造材料}：硅片、特种气体、光掩膜版、光刻胶及配套试剂、CMP抛光材料、靶材、湿化学品等。价值占比约70--75\%。受益标的包括沪硅产业（硅片）、安集科技（抛光液）、江丰电子（靶材）、华特气体（特种气体）、南大光电（光刻胶）等。
\item \textbf{封装材料}：引线框架、封装基板、键合丝、塑封料等。国产化率相对较高，竞争格局更分散。
\end{itemize}

从投资优先级看，我们推荐\textbf{设备优于材料，前道优于后道，关键零部件优于整机}。核心逻辑是：设备的技术壁垒更高、国产替代紧迫性更强、业绩兑现节奏更快；前道设备价值量大、国产化率低，替代空间广阔；核心零部件是产业链最薄弱的环节，一旦突破弹性最大。

\subsection{风险因素与下行空间}

\textbf{主要风险因素：}

\begin{enumerate}[leftmargin=2em]
\item \textbf{下游需求疲软风险}：如果全球经济复苏不及预期，消费电子、服务器等下游需求疲软，可能导致晶圆厂削减资本开支，设备采购放缓。下行空间：若晶圆厂资本开支下调20\%，设备公司订单可能下降15--25\%，对应业绩预期下调和估值承压。

\item \textbf{海外技术封锁升级风险}：如果海外出口管制进一步升级，限制更多设备品类、技术和人才，可能对国内设备厂商的研发和生产造成影响。尤其是关键零部件（如高端射频电源、精密运动部件）的进口受限，可能影响设备交付。

\item \textbf{行业竞争加剧风险}：随着国产替代进程推进，越来越多的企业进入半导体设备材料领域，可能导致中低端市场竞争加剧，价格战压力上升。部分设备品类可能出现“内卷”，影响行业整体盈利能力。
\end{enumerate}

\textbf{下行空间评估}：当前半导体设备头部公司2026E PE中位数约35--40倍，处于近三年40--50\%分位，估值已从2025年的高点回落较多。半导体材料头部公司2026E PE中位数约30--35倍，处于近三年30--40\%分位。若行业景气度显著下行，设备公司估值可能回落至25--30倍区间，材料公司回落至20--25倍区间，对应回调空间约20--30\%。但考虑到国产替代的长期逻辑和政策支撑，深度回调的空间有限，反而可能提供更好的布局时机。

\subsection{关键标的点评}

\textbf{中微公司（688012）：刻蚀设备龙头，技术实力国际一流}

公司是国内刻蚀设备的绝对龙头，技术实力达到国际先进水平，产品覆盖CCP和ICP两大刻蚀技术路线，已进入全球主流晶圆厂供应链。公司刻蚀设备在5nm及以下先进制程的验证持续推进，是国产设备中最有希望率先在先进制程实现大规模量产的公司之一。预计2026--2028年净利润复合增速约30\%，当前对应2026E PE约35倍。

\textbf{北方华创（002371）：综合设备龙头，平台化优势显著}

公司是国内产品线最全的半导体设备厂商，覆盖刻蚀、薄膜沉积、氧化扩散、清洗等多个品类，下游覆盖逻辑芯片、存储芯片、功率半导体、第三代半导体等多个领域。平台化布局使公司能够更好地应对单一品类的周期波动，持续成长性强。预计2026--2028年净利润复合增速约35\%，当前对应2026E PE约38倍。

\textbf{拓荆科技（688072）：薄膜沉积设备领军者，PECVD龙头}

公司是国内PECVD设备的龙头企业，技术实力国内领先，已实现14nm制程设备的验证和批量交付。ALD、SACVD等新产品快速推进，产品矩阵不断丰富。薄膜沉积是半导体设备中价值量最大的品类之一，公司受益于国产替代的弹性大。预计2026--2028年净利润复合增速约40\%，当前对应2026E PE约42倍。

\textbf{沪硅产业（688126）：硅片国产化核心平台，12英寸硅片放量}

公司是国内规模最大的半导体硅片企业，产品覆盖8英寸和12英寸硅片，是国内少数具备12英寸硅片量产能力的厂商。12英寸硅片出货量快速增长，技术节点持续突破，已进入多家主流晶圆厂供应链。受益于国内晶圆厂扩产和硅片国产替代双重逻辑。预计2026--2028年营收复合增速约30\%，2026年有望实现扭亏为盈。

\textbf{安集科技（688019）：CMP抛光液龙头，先进制程突破加速}

公司是国内CMP抛光液的龙头企业，产品在逻辑芯片和存储芯片领域均实现批量供货。公司在先进制程抛光液的研发和验证进展顺利，部分产品已通过14nm以下制程验证。抛光液是半导体材料中技术壁垒较高的品类，一旦通过认证客户粘性极强。预计2026--2028年净利润复合增速约30\%，当前对应2026E PE约32倍。

\begin{keyinsight}[本章小结]
智驾链零部件和半导体设备材料作为中周期拐点型板块的代表，均处于产业趋势明确、市场认知不充分的阶段。智驾链零部件的核心逻辑是\textbf{L3/L4渗透率拐点+国产替代加速+单车价值量提升}三重共振，半导体设备材料的核心逻辑是\textbf{国产替代深化+周期复苏共振+政策持续加持}。两个板块当前估值均处于合理区间，随着催化事件逐步落地，有望迎来业绩与估值的双重提升。
\end{keyinsight}
